\cotefile{91}
\scanpage{47}
pour que $\underline F_1$ ait les \unsure{vertus|conditions} dites. \deleted{\unsure{Il faut}}

$1^\circ)$ $\underline F_1$ doit être stable par multiplication par $\underline B$,
i.e. \deleted{\unsure{les|Soient}} les composantes de $b_i c_\lambda$ suivant
les $d_{\mu}$, soient $\varphi_{i\lambda\mu}\in A_{\unsure{\Phi|\varphi}}$, doivent être
\uline{nulles}, \deleted{\unsure{Il}} d'où des \uline{équations}
\[
\varphi_{i\lambda\mu}=0\qquad
\left[\begin{aligned}
1&\leq i\leq \unsure{\nu|v}\\
1&\leq\lambda\leq(n-k)\nu\\
1&\leq\mu\leq k\nu
\end{aligned}\right]
\]
% slanted
\margin{Soit $G_1$ le sous schéma de $G$ défini par \unsure{ces|ses} équations. On est ramené \unsure{à un cas}}

$2^\circ)$ Le B-module quotient \deleted{$\unsure{B|F}/\unsure{B_1|F_1}$} $F/F_1=M$
est projectif de rang \unsure{en|$n$|$k$}. Il a une
base sur $A$ de la forme $(m_i)_{1\leq i\leq k\nu}$,
Soient $\xi=(\xi_1\mathbin{-}\xi_k)$ \deleted{des} \inserted{un \unsure{syst.} de $k$} \unsure{élts|éléments} de $M$. Pour qu'ils
forment une base sur $B$, il f. et s. que
$B^k\xrightarrow{u_\xi}M$ défini soit bijectif. Or dans $B^k$ et $M$
on a des bases, donc on peut former
$\det(u_\xi)$, la condition est que $\det(u_\xi)$ soit
inversible. \deleted{\unsure{Donc pour un|Donc pour que}} Pour que
en un $\mathrm{p}^{\mathrm t}$ $x\in\operatorname{Spec}(A)$, $M_x$ soit libre sur $B_x$,
il faut et il suffit qu'il existe
$\xi\in(M_x)^k$ tel que $\det(u_\xi)$ soit inversible
dans $\deleted{\unsure{B}}A_x$. On voit facilement que
% slanted
\margin{pour éviter un \unsure{$\forall$|$\mathrm V^{\mathrm t}$}, on \unsure{pourrait expliciter des} conditions \unsure{nécessaires et suffisantes en un}}
l'on peut aussi prendre \unsure{les} $\xi\in M^k$. Donc
pour \deleted{\unsure{tel|tout}} \deleted{$\unsure{\xi\in M}$} $\xi\in M^k$, soit $U_\xi$ l'ouvert
\deleted{\unsure{en}} des non-zéros de $\det(u_\xi)$. Soit $U=\bigcup U_\xi$. C'est
% archive sheet number: 46
